\subsection{Black-hole/wormhole--scattering topological equivalence ladder: certificates and registry (audit)}
\label{app:equivalence_ladder_certificates}

\paragraph{Scope and status.}
\AuditTag This appendix upgrades the ``black-hole/wormhole--scattering topological equivalence ladder'' narrative (Remark~\ref{rem:bh_wh_scattering_topological_equivalence_ladder})
to an auditable \emph{certificate registry}.
It does not assert theorem-level identifications; it records what is claimed, what is measured, and which gates/failure points are invoked.
It also serves as a compact interface for cross-module dictionary alignment, in particular toward the leakage-kernel vocabulary (Appendix~\ref{app:leakage_kernel})
and the graph-zeta/determinant viewpoint (Appendix~\ref{app:graph_zeta_holonomy}).

\paragraph{Certificate format (registry schema).}
\AuditTag Each ladder step is recorded as a certificate tuple:
\[
(\mathrm{Trigger},\ \mathrm{Observable},\ \mathrm{Gate},\ \mathrm{Fallback})
\]
where ``Gate'' must be a named hard gate or minimal failure point instance (Appendix~\ref{app:minimal_failure_point_templates}),
and ``Fallback'' must identify the minimal surviving interface claim if the gate fails.
This keeps the ladder compatible with the manuscript's layer discipline and failure-point normalization.

\paragraph{Registry table (generated).}
\AuditTag The registry rows below are generated by \texttt{scripts/exp\_equivalence\_ladder\_registry.py}.
They point to the concrete toys (Appendix~\ref{app:scattering_bh_toy_equivalence_audits} and Appendix~\ref{app:leakage_kernel}),
to the scattering delay dictionary (Appendix~\ref{app:scattering_haag_ruelle_lsz_interface} and Appendix~\ref{app:time_mass_delay}),
and to the operator bookkeeping dictionary (Appendix~\ref{app:operator_mother_space}).

\begin{table}[t]
\centering
\scriptsize
\setlength{\tabcolsep}{6pt}
\renewcommand{\arraystretch}{1.15}
\begin{tabular}{l l l l l l}
\toprule
id & step & carrier & observables & gates/failure & pointers \\
\midrule
C1 & BH-like saturation $\leftrightarrow$ delay-driven backlog & interface/audit (queue record) & frac\\_q$>q\_\star$, max\\_q; $\overline{\tau}\_{\mathrm{WS}}$ & (S1) unitarity window; declared band/grid & App.~\ref{app:scattering_bh_toy_equivalence_audits}, Tab.~\ref{tab:scattering_bh_queue_equivalence_toy} \\
C2 & leak/evaporation $\leftrightarrow$ decay exit (width/lifetime proxy) & interface/audit (finite families) & $P\_f(t)$, $\Gamma\_f$, $\tau\_f$; exit-channel tokens & finite-family discipline; CAP tie-break & App.~\ref{app:leakage_kernel}, Def.~\ref{def:survival_kernel_family} \\
C3 & scattering delay $\leftrightarrow$ phase/logdet bookkeeping & interface (dictionary) & $\tau\_{\mathrm{WS}}(\omega)=-\iu\,\dd(\log\det S)/\dd\omega$ & (S1) (approx.) unitarity + differentiability on band & App.~\ref{app:scattering_haag_ruelle_lsz_interface}, Lem.~\ref{lem:wigner_smith_trace_logdet} \\
C4 & wormhole-like shortcut $\leftrightarrow$ finite-rank update (operator dictionary) & operator dictionary (resolvent/det) & logdet increment; resolvent-trace delay proxy & (B5)/(IC13) HTF-lite optional; otherwise operator-only baseline & App.~\ref{app:operator_mother_space} (Lem.~\ref{lem:rank_one_det_sherman_morrison}--\ref{lem:finite_rank_det_woodbury}); App.~\ref{app:equivalence_ladder_certificates}, Tab.~\ref{tab:wormhole_finite_rank_delay_logdet_audit} \\
C6 & added edge $\leftrightarrow$ rank-2 update $\leftrightarrow$ det packaging & top/operator (graph det) & $D(z)=\det(I-zA)$ ratio; $\Delta\log|D|$ & finite graph toy; resolvent domain (invertibility) & App.~\ref{app:equivalence_ladder_certificates}, Tab.~\ref{tab:added_edge_det_update_audit} \\
C7 & det phase $\leftrightarrow$ scalar unitary $S(\omega)$ $\leftrightarrow$ delay proxy & frequency-domain packaging (audit) & $\phi(\omega)=\arg D(r\e^{\iu\omega})$; $\tau=\dd\phi/\dd\omega$ & declared (r, omega-grid); phase unwrapping; finite-difference stability & App.~\ref{app:equivalence_ladder_certificates}, Tab.~\ref{tab:det_phase_delay_proxy_audit} \\
C8 & delay proxy trace identity: $\tau=\Im(-\iu z\,\Tr(A(I-zA)^{-1}))$ & operator/determinant (audit identity) & $\tau\_{\mathrm{fd}}$ vs $\tau\_{\mathrm{tr}}$; err bounds & resolvent domain (invertibility); numerical stability & App.~\ref{app:equivalence_ladder_certificates}, Tab.~\ref{tab:det_phase_delay_trace_identity_audit} \\
C9 & unitary $S(\omega)$ proxy: $\tfrac12\Tr Q \leftrightarrow \tau$ (WS dictionary closure) & scattering dictionary (audit toy) & $Q=-\iu S^\dagger S'$; $\tfrac12\Tr Q$ vs $\tau\_{\mathrm{tr}}$ & (S1) is trivial in toy (exact unitary); gate away from det zeros & App.~\ref{app:equivalence_ladder_certificates}, Tab.~\ref{tab:det_to_smatrix_ws_delay_audit} \\
C10 & resonance peak $\leftrightarrow$ linewidth $\Gamma$ $\leftrightarrow$ lifetime proxy & delay/linewidth dictionary (audit toy) & $\tau\_{\max}$; $\Gamma\_{\mathrm{FWHM}}$; $\Gamma\_\tau=4/\tau\_{\max}$ & finite window + grid; peak selection; FWHM existence & App.~\ref{app:equivalence_ladder_certificates}, Tab.~\ref{tab:det_delay_linewidth_audit} \\
C11 & linewidth $\Gamma \Rightarrow$ survival kernel $P(t)=\exp(-\Gamma t)$ (leakage dictionary closure) & decay/evaporation dictionary (audit toy) & $\tau=1/\Gamma$; $P(t)$ at declared times & rate proxy must be nonnegative; declared time grid & App.~\ref{app:equivalence_ladder_certificates}, Tab.~\ref{tab:linewidth_survival_kernel_audit} \\
C12 & finite-family survival kernels: CAP selection and rate proxy $\Gamma_f$ & finite-family discipline (audit) & $\Gamma\_f$; $\tau\_f$; $P\_f(t)$; CAP-selected row & explicit finite family + deterministic tie-break & App.~\ref{app:equivalence_ladder_certificates}, Tab.~\ref{tab:survival_finite_family_audit} \\
C13 & Ihara/Hashimoto/Bass determinant packaging $\leftrightarrow$ added-edge update ratio & graph zeta / det packaging (audit toy) & $\det(I-uB)$; Bass determinant; edge/base ratio & degree>=2 (toy uses cycle); declared u-grid & App.~\ref{app:equivalence_ladder_certificates}, Tab.~\ref{tab:ihara_hashimoto_added_edge_audit} \\
C14 & dimension-changing update: $\det(I-uB_{\mathrm{new}})/\det(I-uB_{\mathrm{old}})=\det(S(u))$ (Schur factor) & graph zeta / block determinant (audit identity) & ratio (direct) vs ratio (Schur); err & det(I-uB\_old) nonzero on grid; explicit block split & App.~\ref{app:equivalence_ladder_certificates}, Tab.~\ref{tab:hashimoto_added_edge_schur_audit} \\
C15 & local interface factorization contract: small det factor (Woodbury/Schur) & operator mother space (audit unification) & det ratio = det(small factor) in both regimes & declared baseline + explicit interface split; no hidden degrees of freedom & App.~\ref{app:operator_mother_space}, Rem.~\ref{rem:local_interface_factorization_contract}; App.~\ref{app:equivalence_ladder_certificates}, Tabs.~\ref{tab:wormhole_finite_rank_delay_logdet_audit},\ref{tab:hashimoto_added_edge_schur_audit} \\
C5 & four-way unification as a contract form & audit thesis (contract) & ledger form + gate checklist & hard gates in full-fusion interface; explicit fallbacks & Prop.~\ref{prop:four_way_unification_contract_form}, Rem.~\ref{rem:four_way_unification_contract_fallback} \\
\bottomrule
\end{tabular}
\caption{Certificate registry for the black-hole/wormhole--scattering topological equivalence ladder (audit-only).}
\label{tab:equivalence_ladder_registry}
\end{table}

\paragraph{Certificate registry (equivalence ladder).} \AuditTag The table records a minimal set of canonical certificates for the BH/WH $\leftrightarrow$ scattering/decay equivalence ladder. Each row is a declared contract: it names the trigger and readout observable, the required gate/failure point instance, and where the corresponding reproducible artifact or dictionary lives in the manuscript.
\paragraph{Discipline.} \AuditTag Rows are not claims of ontological identity. They are audit entries: if the named gate fails, the step must revert to its declared fallback (typically a weaker interface comparison or a pure bookkeeping statement).%

\paragraph{LF1/C15: local interface factorization templates (audit gate).}
\AuditTag The ladder step that upgrades a shortcut/update narrative to global determinant/logdet/phase-delay bookkeeping is governed by the
\emph{local interface factorization} gate (LF1) (Appendix~\ref{app:minimal_failure_point_templates}).
Operationally, the upgrade is admissible only if one exhibits (i) an explicit baseline/counterfactual split and (ii) an explicit \emph{small} determinant factor
that isolates the effect of the local interface modification.
There are two canonical regimes:
\begin{itemize}
\item \textbf{Same-dimension finite-rank update (Woodbury template).}
Let $A$ be a finite operator on a fixed space and $D(z):=\det(I-zA)$.
For a rank-$k$ update $A\mapsto A+U V^\top$ (with $U,V$ having $k$ columns), one has the small-factor ratio
\[
\frac{\det(I-z(A+UV^\top))}{\det(I-zA)}
=\det\!\Bigl(I_k - z\,V^\top (I-zA)^{-1}U\Bigr),
\]
so the bookkeeping delta is forced to factor through a $k\times k$ interface matrix (Appendix~\ref{app:operator_mother_space},
Lemma~\ref{lem:finite_rank_det_woodbury}).
\item \textbf{Dimension-changing update (Schur template).}
When the updated object lives on an enlarged space, write the block split (old | new) so that
\[
M_{\mathrm{new}}=
\begin{pmatrix}
M_{\mathrm{old}} & B\\
C & D
\end{pmatrix},
\qquad
S:=D-CM_{\mathrm{old}}^{-1}B,
\]
then $\det(M_{\mathrm{new}})=\det(M_{\mathrm{old}})\det(S)$.
Thus the ratio is again a small determinant factor $\det(S)$, with size equal to the number of added degrees of freedom (Appendix~\ref{app:operator_mother_space},
Remark~\ref{rem:local_interface_factorization_contract}).
\end{itemize}
If the baseline or the small-factor split is not provided, the step must revert to its declared fallback in the registry
(typically a pure pointer-jump/graph-update statement without determinant/delay promotion).

\paragraph{Finite-rank update $\Rightarrow$ delay/logdet bookkeeping (generated audit).}
\AuditTag To make the ``wormhole-like shortcut $\leftrightarrow$ finite-rank update'' step measurable in a minimal toy,
we provide a deterministic audit artifact that compares a baseline operator to its finite-rank updated version and reports
logdet-style increments together with a resolvent/determinant-derived delay proxy.
Rows and summaries are generated by \texttt{scripts/exp\_wormhole\_finite\_rank\_delay\_logdet\_audit.py}.

\begin{table}[t]
\centering
\scriptsize
\setlength{\tabcolsep}{6pt}
\renewcommand{\arraystretch}{1.15}
\begin{tabular}{r r r r r r r r r}
\toprule
$\alpha$ & cost & $r$ & $\log|\det|_0$ & $\log|\det|_1$ & $\Delta\log|\det|$ & $\tau_0(r)$ & $\tau_1(r)$ & $\Delta\tau$ \\
\midrule
0.000 & 0.000 & 0.20 & -0.143010 & -0.143010 & 0.000000 & 0.730589 & 0.730589 & 0.000000 \\
0.000 & 0.000 & 0.40 & -0.292447 & -0.292447 & 0.000000 & 0.764357 & 0.764357 & 0.000000 \\
0.000 & 0.000 & 0.60 & -0.449003 & -0.449003 & 0.000000 & 0.801874 & 0.801874 & 0.000000 \\
0.000 & 0.000 & 0.80 & -0.613495 & -0.613495 & 0.000000 & 0.843862 & 0.843862 & 0.000000 \\
0.100 & 0.100 & 0.20 & -0.143010 & -0.151752 & -0.008742 & 0.730589 & 0.776895 & 0.046305 \\
0.100 & 0.100 & 0.40 & -0.292447 & -0.311031 & -0.018584 & 0.764357 & 0.816649 & 0.052292 \\
0.100 & 0.100 & 0.60 & -0.449003 & -0.478739 & -0.029737 & 0.801874 & 0.861331 & 0.059457 \\
0.100 & 0.100 & 0.80 & -0.613495 & -0.655962 & -0.042468 & 0.843862 & 0.912001 & 0.068139 \\
0.200 & 0.200 & 0.20 & -0.143010 & -0.160571 & -0.017560 & 0.730589 & 0.824020 & 0.093431 \\
0.200 & 0.200 & 0.40 & -0.292447 & -0.329967 & -0.037520 & 0.764357 & 0.870941 & 0.106583 \\
0.200 & 0.200 & 0.60 & -0.449003 & -0.509388 & -0.060385 & 0.801874 & 0.924489 & 0.122615 \\
0.200 & 0.200 & 0.80 & -0.613495 & -0.700314 & -0.086819 & 0.843862 & 0.986320 & 0.142458 \\
0.300 & 0.300 & 0.20 & -0.143010 & -0.169468 & -0.026458 & 0.730589 & 0.871988 & 0.141399 \\
0.300 & 0.300 & 0.40 & -0.292447 & -0.349269 & -0.056822 & 0.764357 & 0.927348 & 0.162991 \\
0.300 & 0.300 & 0.60 & -0.449003 & -0.541005 & -0.092003 & 0.801874 & 0.991705 & 0.189831 \\
0.300 & 0.300 & 0.80 & -0.613495 & -0.746724 & -0.133230 & 0.843862 & 1.067700 & 0.223838 \\
\bottomrule
\end{tabular}
\caption{Toy audit for a rank-one pointer-jump update: determinant bookkeeping and resolvent-derived delay proxy increments.}
\label{tab:wormhole_finite_rank_delay_logdet_audit}
\end{table}

\paragraph{Finite-rank update induces logdet and delay-proxy increments (toy).} \AuditTag We compare a baseline finite operator $F$ to a rank-one updated operator $F+\Delta$ (``wormhole-on''), where the update amplitude $\alpha$ is treated as an explicit ledger-cost proxy. On an Abel-like parameter grid $r\in(0,1)$ we report: (i) a logdet bookkeeping proxy $\log|\det(I-rF)|$ and its increment, and (ii) a determinant/resolvent-derived delay proxy $\tau(r)=\mathrm{Tr}\,F(I-rF)^{-1}$ and its increment (via $\frac{\dd}{\dd r}\log\det(I-rF)=-\tau(r)$).
\paragraph{Deterministic witness on the declared grid.} \AuditTag Over the fixed grid (r=[0.2, 0.4, 0.6, 0.8], alpha=[0.0, 0.1, 0.2, 0.3]), the largest observed increments are $\Delta\tau\approx 0.223838$ at $(\alpha,r)=(0.300,0.80)$ and $\Delta\log|\det|\approx -0.133230$ at $(\alpha,r)=(0.300,0.80)$. The full row table is provided in \texttt{sections/generated/wormhole\_finite\_rank\_delay\_logdet\_rows.tex}.%

\paragraph{Added edge $\Rightarrow$ finite-rank update $\Rightarrow$ determinant packaging (generated audit).}
\AuditTag The ladder step ``added shortcut edge'' can be audited directly in a fixed finite graph toy:
adding one undirected edge is a rank-2 update of the adjacency operator, and the determinant package
$D(z)=\det(I-zA)$ changes by a finite-rank determinant factor computable from a $2\times2$ matrix (Woodbury).
Rows and summary are generated by \texttt{scripts/exp\_added\_edge\_det\_update\_audit.py}.

\begin{table}[t]
\centering
\scriptsize
\setlength{\tabcolsep}{6pt}
\renewcommand{\arraystretch}{1.15}
\begin{tabular}{r r r r r r r r r}
\toprule
$n$ & edge & $z$ & $\det_0$ & $\det_1$ & ratio (direct) & ratio (Woodbury) & $\Delta\log|\det|$ & err \\
\midrule
6 & (0,3) & 0.050 & 0.987537 & 0.985038 & 0.997468 & 0.997468 & -0.002535 & 0.000000 \\
6 & (0,3) & 0.100 & 0.950599 & 0.940600 & 0.989481 & 0.989481 & -0.010574 & 0.000000 \\
6 & (0,3) & 0.150 & 0.890526 & 0.868038 & 0.974747 & 0.974747 & -0.025578 & 0.000000 \\
6 & (0,3) & 0.200 & 0.809536 & 0.769600 & 0.950668 & 0.950668 & -0.050590 & 0.000000 \\
6 & (0,3) & 0.250 & 0.710693 & 0.648438 & 0.912401 & 0.912401 & -0.091675 & 0.000000 \\
\bottomrule
\end{tabular}
\caption{Added-edge toy: determinant package ratio computed directly and via a $2\times2$ Woodbury factor.}
\label{tab:added_edge_det_update_audit}
\end{table}

\paragraph{Added edge $\leftrightarrow$ finite-rank update (det packaging toy).} \AuditTag On a fixed 6-node path graph adjacency $A_0$ we add a shortcut edge $(0,3)$, obtaining $A_1=A_0+\Delta$ with a rank-2 update $\Delta=e_0e_3^\top+e_3e_0^\top$. We compare determinant packages $D(z)=\det(I-zA)$ and report the ratio $D_1(z)/D_0(z)$ and the logabs increment $\Delta\log|D|$. The ratio is also computed via the finite-rank determinant lemma (Woodbury) as a $2\times2$ determinant built from $(I-zA_0)^{-1}$.
\paragraph{Deterministic agreement on the declared grid.} \AuditTag On z-grid [0.05, 0.1, 0.15, 0.2, 0.25], the maximum absolute discrepancy between the direct ratio and the Woodbury ratio is 0.000000 (at z=0.100).%

\paragraph{Determinant phase $\Rightarrow$ delay proxy (frequency-domain packaging; generated audit).}
\AuditTag To align determinant packaging with the ``phase--delay--logdet'' dictionary used on the scattering side,
we define a scalar unitary proxy on a declared frequency grid by sampling $D(r\e^{\iu\omega})=\det(I-r\e^{\iu\omega}A)$ at fixed $r<1$ and setting
$S(\omega)=\exp(\iu\,\arg D(r\e^{\iu\omega}))$.
The associated delay proxy is the phase derivative $\tau(\omega)=\dd(\arg D)/\dd\omega$ (finite-difference audited).
Rows and summary are generated by \texttt{scripts/exp\_det\_phase\_delay\_scattering\_proxy\_audit.py}.

\begin{table}[t]
\centering
\scriptsize
\setlength{\tabcolsep}{6pt}
\renewcommand{\arraystretch}{1.15}
\begin{tabular}{r r r r r r r r r r}
\toprule
$n$ & edge & $r$ & $\omega$ & $\phi_0$ & $\phi_1$ & $\Delta\phi$ & $\tau_0$ & $\tau_1$ & $\Delta\tau$ \\
\midrule
6 & (0,3) & 0.350 & 0.000 & 0.000000 & 0.000000 & 0.000000 & -1.802235 & -2.969772 & -1.167536 \\\\
6 & (0,3) & 0.350 & 0.050 & -0.091566 & -0.154236 & -0.062671 & -1.802235 & -2.969772 & -1.167536 \\\\
6 & (0,3) & 0.350 & 0.100 & -0.180224 & -0.296977 & -0.116754 & -1.718414 & -2.661701 & -0.943287 \\\\
6 & (0,3) & 0.350 & 0.150 & -0.263407 & -0.420406 & -0.156999 & -1.589063 & -2.245206 & -0.656144 \\\\
6 & (0,3) & 0.350 & 0.200 & -0.339130 & -0.521498 & -0.182368 & -1.426724 & -1.804027 & -0.377303 \\\\
6 & (0,3) & 0.350 & 0.250 & -0.406079 & -0.600809 & -0.194730 & -1.244458 & -1.391901 & -0.147443 \\\\
6 & (0,3) & 0.350 & 0.300 & -0.463576 & -0.660688 & -0.197112 & -1.053729 & -1.031722 & 0.022007 \\\\
6 & (0,3) & 0.350 & 0.350 & -0.511452 & -0.703981 & -0.192529 & -0.863366 & -0.726907 & 0.136459 \\\\
6 & (0,3) & 0.350 & 0.400 & -0.549912 & -0.733379 & -0.183466 & -0.679441 & -0.471999 & 0.207442 \\\\
6 & (0,3) & 0.350 & 0.450 & -0.579396 & -0.751181 & -0.171785 & -0.505677 & -0.258928 & 0.246749 \\\\
6 & (0,3) & 0.350 & 0.500 & -0.600480 & -0.759271 & -0.158792 & -0.344028 & -0.079849 & 0.264179 \\\\
6 & (0,3) & 0.350 & 0.550 & -0.613799 & -0.759166 & -0.145367 & -0.195239 & 0.071882 & 0.267121 \\\\
6 & (0,3) & 0.350 & 0.600 & -0.620004 & -0.752083 & -0.132079 & -0.059286 & 0.201565 & 0.260851 \\\\
6 & (0,3) & 0.350 & 0.650 & -0.619728 & -0.739009 & -0.119282 & 0.064306 & 0.313321 & 0.249015 \\\\
6 & (0,3) & 0.350 & 0.700 & -0.613573 & -0.720751 & -0.107178 & 0.176247 & 0.410324 & 0.234077 \\\\
6 & (0,3) & 0.350 & 0.750 & -0.602103 & -0.697977 & -0.095874 & 0.277349 & 0.495019 & 0.217670 \\\\
6 & (0,3) & 0.350 & 0.800 & -0.585838 & -0.671249 & -0.085411 & 0.368446 & 0.569306 & 0.200860 \\\\
6 & (0,3) & 0.350 & 0.850 & -0.565258 & -0.641047 & -0.075788 & 0.450343 & 0.634671 & 0.184327 \\\\
6 & (0,3) & 0.350 & 0.900 & -0.540804 & -0.607782 & -0.066978 & 0.523791 & 0.692291 & 0.168499 \\\\
6 & (0,3) & 0.350 & 0.950 & -0.512879 & -0.571817 & -0.058938 & 0.589477 & 0.743110 & 0.153632 \\\\
6 & (0,3) & 0.350 & 1.000 & -0.481856 & -0.533471 & -0.051615 & 0.648019 & 0.787891 & 0.139873 \\\\
6 & (0,3) & 0.350 & 1.050 & -0.448077 & -0.493028 & -0.044951 & 0.699963 & 0.827259 & 0.127296 \\\\
6 & (0,3) & 0.350 & 1.100 & -0.411860 & -0.450745 & -0.038885 & 0.745794 & 0.861729 & 0.115935 \\\\
6 & (0,3) & 0.350 & 1.150 & -0.373498 & -0.406855 & -0.033357 & 0.785930 & 0.891723 & 0.105793 \\\\
6 & (0,3) & 0.350 & 1.200 & -0.333267 & -0.361573 & -0.028306 & 0.820736 & 0.917597 & 0.096861 \\\\
6 & (0,3) & 0.350 & 1.250 & -0.291424 & -0.315096 & -0.023671 & 0.850522 & 0.939641 & 0.089119 \\\\
6 & (0,3) & 0.350 & 1.300 & -0.248215 & -0.267609 & -0.019394 & 0.875549 & 0.958095 & 0.082546 \\\\
6 & (0,3) & 0.350 & 1.350 & -0.203870 & -0.219286 & -0.015417 & 0.896035 & 0.973155 & 0.077121 \\\\
6 & (0,3) & 0.350 & 1.400 & -0.158611 & -0.170293 & -0.011682 & 0.912154 & 0.984978 & 0.072824 \\\\
6 & (0,3) & 0.350 & 1.450 & -0.112654 & -0.120789 & -0.008134 & 0.924044 & 0.993683 & 0.069639 \\\\
6 & (0,3) & 0.350 & 1.500 & -0.066207 & -0.070925 & -0.004718 & 0.931804 & 0.999357 & 0.067553 \\\\
6 & (0,3) & 0.350 & 1.550 & -0.019474 & -0.020853 & -0.001379 & 0.935499 & 1.002057 & 0.066558 \\\\
6 & (0,3) & 0.350 & 1.600 & 0.027343 & 0.029281 & 0.001938 & 0.935161 & 1.001809 & 0.066649 \\\\
6 & (0,3) & 0.350 & 1.650 & 0.074042 & 0.079328 & 0.005286 & 0.930785 & 0.998612 & 0.067827 \\\\
6 & (0,3) & 0.350 & 1.700 & 0.120422 & 0.129142 & 0.008720 & 0.922336 & 0.992433 & 0.070097 \\\\
6 & (0,3) & 0.350 & 1.750 & 0.166276 & 0.178571 & 0.012296 & 0.909743 & 0.983211 & 0.073468 \\\\
6 & (0,3) & 0.350 & 1.800 & 0.211396 & 0.227463 & 0.016067 & 0.892900 & 0.970853 & 0.077954 \\\\
6 & (0,3) & 0.350 & 1.850 & 0.255566 & 0.275657 & 0.020091 & 0.871664 & 0.955234 & 0.083570 \\\\
6 & (0,3) & 0.350 & 1.900 & 0.298562 & 0.322986 & 0.024424 & 0.845853 & 0.936191 & 0.090338 \\\\
6 & (0,3) & 0.350 & 1.950 & 0.340151 & 0.369276 & 0.029125 & 0.815243 & 0.913521 & 0.098279 \\\\
6 & (0,3) & 0.350 & 2.000 & 0.380087 & 0.414338 & 0.034252 & 0.815243 & 0.913521 & 0.098279 \\\\
\bottomrule
\end{tabular}
\caption{Frequency-domain proxy: determinant phase and its delay proxy under an added-edge update.}
\label{tab:det_phase_delay_proxy_audit}
\end{table}

\paragraph{Determinant packaging $\Rightarrow$ scalar phase and delay proxy (toy).} \AuditTag For a fixed finite graph adjacency $A$ we define the determinant package $D(z)=\det(I-zA)$. On a fixed Abel-like radius $r<1$ we sample $z=r\e^{\iu\omega}$ and take the phase $\phi(\omega)=\arg D(r\e^{\iu\omega})$. This induces a scalar unitary proxy $S(\omega)=\exp(\iu\phi(\omega))$, whose delay proxy is $\tau(\omega)=\dd\phi/\dd\omega$. We compare base vs added-edge update and report the increments $\Delta\phi$ and $\Delta\tau$ on a declared grid.
\paragraph{Deterministic grid witness.} \AuditTag With n=6, added edge (0,3), r=0.350, omega-grid size k=41 with step 0.05, the maximal observed magnitudes are max|Δφ|≈0.197112 and max|Δτ|≈1.167536.%

\paragraph{Trace identity: $\tau(\omega)$ from $\Tr(A(I-zA)^{-1})$ (generated audit).}
\AuditTag The finite-difference delay proxy above can be audited against a closed trace identity.
For $D(z)=\det(I-zA)$ one has $\frac{\dd}{\dd z}\log D(z)=-\Tr(A(I-zA)^{-1})$, hence on $z=r\e^{\iu\omega}$,
\[
\tau(\omega)=\frac{\dd}{\dd\omega}\arg D(r\e^{\iu\omega})
=\Im\frac{\dd}{\dd\omega}\log D(r\e^{\iu\omega})
=\Im\bigl(-\iu z\,\Tr(A(I-zA)^{-1})\bigr).
\]
The table below compares the phase-difference estimate (fd) to the resolvent-trace formula (tr) for both wormhole-off/on (base vs added edge) and for $\Delta\tau$.
Rows and summary are generated by \texttt{scripts/exp\_det\_phase\_delay\_scattering\_proxy\_audit.py}.

\begin{table}[t]
\centering
\scriptsize
\setlength{\tabcolsep}{6pt}
\renewcommand{\arraystretch}{1.15}
\begin{tabular}{r r r r r r r r r r}
\toprule
$\omega$ & $\tau_{0,\mathrm{fd}}$ & $\tau_{0,\mathrm{tr}}$ & err$_0$
& $\tau_{1,\mathrm{fd}}$ & $\tau_{1,\mathrm{tr}}$ & err$_1$
& $\Delta\tau_{\mathrm{fd}}$ & $\Delta\tau_{\mathrm{tr}}$ & err$_\Delta$ \\
\midrule
0.000 & -1.802235 & -1.841247 & 0.039011 & -2.969772 & -3.126008 & 0.156236 & -1.167536 & -1.284761 & 0.117225 \\\\
0.050 & -1.802235 & -1.811565 & 0.009330 & -2.969772 & -3.003742 & 0.033970 & -1.167536 & -1.192177 & 0.024640 \\\\
0.100 & -1.718414 & -1.726115 & 0.007701 & -2.661701 & -2.679503 & 0.017803 & -0.943287 & -0.953389 & 0.010102 \\\\
0.150 & -1.589063 & -1.594572 & 0.005509 & -2.245206 & -2.248081 & 0.002875 & -0.656144 & -0.653509 & 0.002635 \\\\
0.200 & -1.426724 & -1.429983 & 0.003259 & -1.804027 & -1.798086 & 0.005941 & -0.377303 & -0.368103 & 0.009200 \\\\
0.250 & -1.244458 & -1.245774 & 0.001316 & -1.391901 & -1.382618 & 0.009283 & -0.147443 & -0.136844 & 0.010599 \\\\
0.300 & -1.053729 & -1.053573 & 0.000156 & -1.031722 & -1.022236 & 0.009486 & 0.022007 & 0.031337 & 0.009330 \\\\
0.350 & -0.863366 & -0.862213 & 0.001153 & -0.726907 & -0.718526 & 0.008381 & 0.136459 & 0.143687 & 0.007228 \\\\
0.400 & -0.679441 & -0.677689 & 0.001752 & -0.471999 & -0.465046 & 0.006953 & 0.207442 & 0.212643 & 0.005201 \\\\
0.450 & -0.505677 & -0.503618 & 0.002059 & -0.258928 & -0.253308 & 0.005620 & 0.246749 & 0.250310 & 0.003561 \\\\
0.500 & -0.344028 & -0.341860 & 0.002168 & -0.079849 & -0.075338 & 0.004512 & 0.264179 & 0.266523 & 0.002344 \\\\
0.550 & -0.195239 & -0.193086 & 0.002154 & 0.071882 & 0.075516 & 0.003635 & 0.267121 & 0.268602 & 0.001481 \\\\
0.600 & -0.059286 & -0.057218 & 0.002068 & 0.201565 & 0.204521 & 0.002957 & 0.260851 & 0.261739 & 0.000889 \\\\
0.650 & 0.064306 & 0.066251 & 0.001945 & 0.313321 & 0.315756 & 0.002435 & 0.249015 & 0.249505 & 0.000490 \\\\
0.700 & 0.176247 & 0.178053 & 0.001807 & 0.410324 & 0.412357 & 0.002033 & 0.234077 & 0.234303 & 0.000226 \\\\
0.750 & 0.277349 & 0.279016 & 0.001667 & 0.495019 & 0.496741 & 0.001721 & 0.217670 & 0.217725 & 0.000055 \\\\
0.800 & 0.368446 & 0.369978 & 0.001532 & 0.569306 & 0.570783 & 0.001477 & 0.200860 & 0.200805 & 0.000055 \\\\
0.850 & 0.450343 & 0.451749 & 0.001406 & 0.634671 & 0.635954 & 0.001283 & 0.184327 & 0.184205 & 0.000123 \\\\
0.900 & 0.523791 & 0.525083 & 0.001291 & 0.692291 & 0.693418 & 0.001128 & 0.168499 & 0.168336 & 0.000164 \\\\
0.950 & 0.589477 & 0.590666 & 0.001188 & 0.743110 & 0.744111 & 0.001002 & 0.153632 & 0.153445 & 0.000187 \\\\
1.000 & 0.648019 & 0.649116 & 0.001097 & 0.787891 & 0.788790 & 0.000899 & 0.139873 & 0.139674 & 0.000199 \\\\
1.050 & 0.699963 & 0.700980 & 0.001017 & 0.827259 & 0.828073 & 0.000814 & 0.127296 & 0.127093 & 0.000203 \\\\
1.100 & 0.745794 & 0.746741 & 0.000947 & 0.861729 & 0.862472 & 0.000743 & 0.115935 & 0.115731 & 0.000204 \\\\
1.150 & 0.785930 & 0.786817 & 0.000887 & 0.891723 & 0.892408 & 0.000685 & 0.105793 & 0.105591 & 0.000202 \\\\
1.200 & 0.820736 & 0.821571 & 0.000835 & 0.917597 & 0.918233 & 0.000637 & 0.096861 & 0.096662 & 0.000199 \\\\
1.250 & 0.850522 & 0.851314 & 0.000792 & 0.939641 & 0.940237 & 0.000597 & 0.089119 & 0.088924 & 0.000195 \\\\
1.300 & 0.875549 & 0.876305 & 0.000756 & 0.958095 & 0.958659 & 0.000564 & 0.082546 & 0.082355 & 0.000191 \\\\
1.350 & 0.896035 & 0.896761 & 0.000726 & 0.973155 & 0.973694 & 0.000538 & 0.077121 & 0.076933 & 0.000188 \\\\
1.400 & 0.912154 & 0.912858 & 0.000704 & 0.984978 & 0.985496 & 0.000519 & 0.072824 & 0.072639 & 0.000185 \\\\
1.450 & 0.924044 & 0.924731 & 0.000687 & 0.993683 & 0.994187 & 0.000504 & 0.069639 & 0.069456 & 0.000183 \\\\
1.500 & 0.931804 & 0.932480 & 0.000676 & 0.999357 & 0.999852 & 0.000495 & 0.067553 & 0.067371 & 0.000182 \\\\
1.550 & 0.935499 & 0.936171 & 0.000671 & 1.002057 & 1.002547 & 0.000490 & 0.066558 & 0.066377 & 0.000181 \\\\
1.600 & 0.935161 & 0.935832 & 0.000672 & 1.001809 & 1.002300 & 0.000491 & 0.066649 & 0.066468 & 0.000181 \\\\
1.650 & 0.930785 & 0.931463 & 0.000678 & 0.998612 & 0.999108 & 0.000496 & 0.067827 & 0.067645 & 0.000182 \\\\
1.700 & 0.922336 & 0.923025 & 0.000689 & 0.992433 & 0.992939 & 0.000506 & 0.070097 & 0.069914 & 0.000183 \\\\
1.750 & 0.909743 & 0.910450 & 0.000707 & 0.983211 & 0.983732 & 0.000522 & 0.073468 & 0.073283 & 0.000186 \\\\
1.800 & 0.892900 & 0.893631 & 0.000731 & 0.970853 & 0.971396 & 0.000542 & 0.077954 & 0.077765 & 0.000188 \\\\
1.850 & 0.871664 & 0.872425 & 0.000761 & 0.955234 & 0.955803 & 0.000569 & 0.083570 & 0.083378 & 0.000192 \\\\
1.900 & 0.845853 & 0.846651 & 0.000798 & 0.936191 & 0.936794 & 0.000603 & 0.090338 & 0.090143 & 0.000195 \\\\
1.950 & 0.815243 & 0.816086 & 0.000843 & 0.913521 & 0.914166 & 0.000644 & 0.098279 & 0.098080 & 0.000199 \\\\
2.000 & 0.815243 & 0.780460 & 0.034783 & 0.913521 & 0.887671 & 0.025850 & 0.098279 & 0.107211 & 0.008933 \\\\
\bottomrule
\end{tabular}
\caption{Trace identity audit for the frequency-domain delay proxy derived from determinant packaging.}
\label{tab:det_phase_delay_trace_identity_audit}
\end{table}

\paragraph{Trace identity audit: phase-derivative delay vs resolvent-trace formula.} \AuditTag For $D(z)=\det(I-zA)$ one has $\frac{\dd}{\dd z}\log D(z)=-\Tr(A(I-zA)^{-1})$. On $z=r\e^{\iu\omega}$ this yields the phase-derivative delay proxy $\tau(\omega)=\Im\frac{\dd}{\dd\omega}\log D(r\e^{\iu\omega})=\Im\bigl(-\iu z\,\Tr(A(I-zA)^{-1})\bigr)$. We compare this trace formula to the finite-difference estimate from the unwrapped phase on the same declared grid.
\paragraph{Deterministic error on the declared grid.} \AuditTag With r=0.350 and domega=0.05, we gate phase-difference auditing by requiring $|D(r\e^{\iu\omega})|\ge 0.001$. Over the remaining points, the maximum absolute discrepancies are $\max|\tau_{0,\mathrm{fd}}-\tau_{0,\mathrm{tr}}|\approx0.039011$, $\max|\tau_{1,\mathrm{fd}}-\tau_{1,\mathrm{tr}}|\approx0.156236$, and $\max|\Delta\tau_{\mathrm{fd}}-\Delta\tau_{\mathrm{tr}}|\approx0.117225$.%

\paragraph{Unitary $S(\omega)$ proxy and Wigner--Smith delay (generated audit).}
\AuditTag To make the scattering-side dictionary object explicit, we form a strict $2\times2$ unitary toy
$S(\omega)=\exp(\iu\,\phi(\omega))I_2$ from the determinant phase $\phi(\omega)=\arg D(r\e^{\iu\omega})$.
Then the Wigner--Smith operator satisfies $Q(\omega)=-\iu S^\dagger \frac{\dd S}{\dd\omega}=(\dd\phi/\dd\omega)\,I_2$ and hence $\tfrac12\Tr Q=\dd\phi/\dd\omega$.
The table below audits that this WS-delay proxy matches the resolvent-trace delay identity on the same declared grid (with the same gate away from determinant zeros).
Rows and summary are generated by \texttt{scripts/exp\_det\_to\_unitary\_smatrix\_ws\_delay\_audit.py}.

\begin{table}[t]
\centering
\scriptsize
\setlength{\tabcolsep}{6pt}
\renewcommand{\arraystretch}{1.15}
\begin{tabular}{r r r r r r r r r r}
\toprule
$\omega$ & $\tau_{0,\mathrm{tr}}$ & $\tau_{0,\mathrm{ws}}$ & err$_0$
& $\tau_{1,\mathrm{tr}}$ & $\tau_{1,\mathrm{ws}}$ & err$_1$
& $\Delta\tau_{\mathrm{tr}}$ & $\Delta\tau_{\mathrm{ws}}$ & err$_\Delta$ \\
\midrule
\input{sections/generated/det_to_smatrix_ws_delay_rows.tex}
\end{tabular}
\caption{Unitary $S(\omega)$ toy: Wigner--Smith delay proxy versus resolvent-trace delay identity.}
\label{tab:det_to_smatrix_ws_delay_audit}
\end{table}

\paragraph{Det phase $\Rightarrow$ unitary $S(\omega)$ and WS delay (toy).} \AuditTag Given a determinant package $D(r\e^{\iu\omega})=\det(I-r\e^{\iu\omega}A)$ we form the unwrapped phase $\phi(\omega)=\arg D(r\e^{\iu\omega})$ and define a strict $2\times2$ unitary toy $S(\omega)=\exp(\iu\phi(\omega))I_2$. Then the Wigner--Smith operator satisfies $Q(\omega)=-\iu S^\dagger \frac{\dd S}{\dd\omega}=(\dd\phi/\dd\omega)\,I_2$, hence $\tfrac12\Tr Q=\dd\phi/\dd\omega$. We compare this WS delay proxy to the resolvent-trace delay identity $\tau_{\mathrm{tr}}(\omega)=\Im(-\iu z\,\Tr(A(I-zA)^{-1}))$ with $z=r\e^{\iu\omega}$.
\paragraph{Gate and deterministic error.} \AuditTag We gate phase-difference auditing by requiring $|D(r\e^{\iu\omega})|\ge 0.001$ and use domega=0.05. Over the gated points (count=41), the maximum absolute discrepancies are $\max|\tau_{0,\mathrm{ws}}-\tau_{0,\mathrm{tr}}|\approx0.039011$, $\max|\tau_{1,\mathrm{ws}}-\tau_{1,\mathrm{tr}}|\approx0.156236$, and $\max|\Delta\tau_{\mathrm{ws}}-\Delta\tau_{\mathrm{tr}}|\approx0.117225$.%

\paragraph{Resonance $\Rightarrow$ linewidth $\Gamma$ $\Rightarrow$ lifetime proxy (generated audit).}
\AuditTag To close the ``delay/linewidth'' segment of the ladder within the determinant packaging carrier,
we detect resonance peaks in the resolvent-trace delay proxy $\tau(\omega)$ and extract a finite-window linewidth proxy by the full width at half maximum (FWHM).
We also report the Breit--Wigner benchmark proxy $\Gamma_{\tau}=4/\tau_{\max}$ and the mismatch between the two linewidth estimates.
Rows and summary are generated by \texttt{scripts/exp\_det\_delay\_resonance\_linewidth\_audit.py}.

\begin{table}[t]
\centering
\scriptsize
\setlength{\tabcolsep}{6pt}
\renewcommand{\arraystretch}{1.15}
\begin{tabular}{l r r r r r r r}
\toprule
case & rank & $r$ & $\omega_0$ & $\tau_{\max}$ & $\Gamma_{\tau}$ & $\Gamma_{\mathrm{FWHM}}$ & mismatch \\
\midrule
base & 1 & 0.350 & 1.570000 & 0.936518 & 4.271140 & 1.420003 & 2.851137 \\\\
base & 2 & 0.350 & 4.710000 & 0.936514 & 4.271159 & 1.419982 & 2.851177 \\\\
edge & 1 & 0.350 & 1.570000 & 1.002801 & 3.988828 & 1.635624 & 2.353204 \\\\
edge & 2 & 0.350 & 4.710000 & 1.002798 & 3.988839 & 1.635633 & 2.353206 \\\\
\bottomrule
\end{tabular}
\caption{Determinant/resolvent delay peaks and linewidth proxies (toy).}
\label{tab:det_delay_linewidth_audit}
\end{table}

\paragraph{Resonance linewidth proxy from determinant/resolvent delay (toy).} \AuditTag Using the resolvent-trace delay proxy $\tau(\omega)=\Im(-\iu z\,\Tr(A(I-zA)^{-1}))$ at $z=r\e^{\iu\omega}$, we detect the highest positive peaks and extract a full width at half maximum (FWHM) as a linewidth proxy $\Gamma_{\mathrm{FWHM}}$. We also report the Breit--Wigner consistency proxy $\Gamma_{\tau}:=4/\tau_{\max}$ and their mismatch. This is an audit-only finite-window certificate aligning det packaging, phase-delay readouts, and linewidth vocabulary.
\paragraph{Determinism.} \AuditTag Parameters are fixed (n=6, added edge (0,3), r=0.350, omega in [0,2pi] with step 0.01). Only the top 3 positive peaks are reported for each case (base vs added-edge).%

\paragraph{$\Gamma\Rightarrow\tau\Rightarrow P(t)$ as a leakage-kernel entry (generated audit).}
\AuditTag To connect the linewidth/return-rate proxy $\Gamma$ to the decay/evaporation dictionary, we interpret $\Gamma$ as an exit-rate proxy and define
an exponential survival kernel $P(t)=\exp(-\Gamma t)$ with lifetime proxy $\tau=1/\Gamma$ (cf.\ Appendix~\ref{app:leakage_kernel}, Definition~\ref{def:survival_kernel_family} and Remark~\ref{rem:exponential_reference}).
Using the determinant/resolvent carrier, we report two width proxies ($\Gamma_\tau=4/\tau_{\max}$ and $\Gamma_{\mathrm{FWHM}}$) and their induced survival values on declared times.
Rows and summary are generated by \texttt{scripts/exp\_linewidth\_to\_survival\_kernel\_audit.py}.

\begin{table}[t]
\centering
\scriptsize
\setlength{\tabcolsep}{6pt}
\renewcommand{\arraystretch}{1.15}
\begin{tabular}{l r r r r r r r r r}
\toprule
case & $r$ & $\tau_{\max}$ & $\Gamma_{\tau}$ & $\Gamma_{\mathrm{FWHM}}$ & $\tau_{\tau}$ & $\tau_{\mathrm{FWHM}}$ & $t$ & $P_{\tau}(t)$ & $P_{\mathrm{FWHM}}(t)$ \\
\midrule
base & 0.350 & 0.936518 & 4.271140 & 1.420003 & 0.234130 & 0.704224 & 0.250 & 0.343769 & 0.701173 \\\\
base & 0.350 & 0.936518 & 4.271140 & 1.420003 & 0.234130 & 0.704224 & 0.500 & 0.118177 & 0.491644 \\\\
base & 0.350 & 0.936518 & 4.271140 & 1.420003 & 0.234130 & 0.704224 & 1.000 & 0.013966 & 0.241713 \\\\
edge & 0.350 & 1.002801 & 3.988828 & 1.635624 & 0.250700 & 0.611387 & 0.250 & 0.368908 & 0.664377 \\\\
edge & 0.350 & 1.002801 & 3.988828 & 1.635624 & 0.250700 & 0.611387 & 0.500 & 0.136093 & 0.441396 \\\\
edge & 0.350 & 1.002801 & 3.988828 & 1.635624 & 0.250700 & 0.611387 & 1.000 & 0.018521 & 0.194831 \\\\
\bottomrule
\end{tabular}
\caption{Linewidth to survival-kernel audit (toy).}
\label{tab:linewidth_survival_kernel_audit}
\end{table}

\paragraph{Linewidth $\Gamma$ to survival kernel (toy).} \AuditTag We map the linewidth vocabulary to the leakage-kernel dictionary by treating a width proxy $\Gamma$ as an exit-rate proxy and defining the exponential survival law $P(t)=\exp(-\Gamma t)$ with lifetime proxy $\tau=1/\Gamma$. Here $\Gamma$ is extracted from the determinant/resolvent delay carrier by two auditable proxies: $\Gamma_\tau=4/\tau_{\max}$ (Breit--Wigner benchmark) and $\Gamma_{\mathrm{FWHM}}$ (half-maximum width). The table reports both along with the implied survival values on declared times.
\paragraph{Determinism.} \AuditTag Parameters are fixed (n=6, added edge (0,3), r=0.350, omega step 0.01, times=[0.25, 0.5, 1.0]).%

\paragraph{Finite-family survival-kernel CAP selection (generated audit).}
\AuditTag To fully match the leakage-kernel discipline (explicit finite family + deterministic tie-break),
we also provide a CAP selection audit over a finite candidate family of survival kernels $P_f(t)$,
reporting the rate proxy $\Gamma_f$, lifetime proxy $\tau_f=1/\Gamma_f$, and the survival values on declared times.
Rows and summary are generated by \texttt{scripts/exp\_linewidth\_survival\_finite\_family\_audit.py}.

\begin{table}[t]
\centering
\scriptsize
\setlength{\tabcolsep}{6pt}
\renewcommand{\arraystretch}{1.15}
\begin{tabular}{l l r r r r r r r r r r r l}
\toprule
case & family & $\gamma$ & $\alpha$ & $\beta$-scale & $\Gamma_f$ & $\tau_f$ & $t_1$ & $P_f(t_1)$ & $t_2$ & $P_f(t_2)$ & $t_3$ & $P_f(t_3)$ & CAP \\
\midrule
base & exp & 1.420003 & 1.000 & 0.000 & 1.420003 & 0.704224 & 0.25 & 0.701173 & 0.50 & 0.491643 & 1.00 & 0.241713 & yes \\
base & stretched & 1.420003 & 1.500 & 0.000 & 1.420003 & 0.704224 & 0.25 & 0.809356 & 0.50 & 0.549768 & 1.00 & 0.184127 &  \\
base & stretched & 1.420003 & 2.000 & 0.000 & 1.420003 & 0.704224 & 0.25 & 0.881592 & 0.50 & 0.604048 & 1.00 & 0.133133 &  \\
base & logquad & 1.420003 & 1.000 & 0.200 & 1.420003 & 0.704224 & 0.25 & 0.683721 & 0.50 & 0.444493 & 1.00 & 0.161494 &  \\
base & logquad & 1.420003 & 1.000 & 0.500 & 1.420003 & 0.704224 & 0.25 & 0.658353 & 0.50 & 0.382108 & 1.00 & 0.088195 &  \\
base & exp & 2.845572 & 1.000 & 0.000 & 2.845572 & 0.351423 & 0.25 & 0.490960 & 0.50 & 0.241042 & 1.00 & 0.058101 &  \\
base & stretched & 2.845572 & 1.500 & 0.000 & 2.845572 & 0.351423 & 0.25 & 0.548802 & 0.50 & 0.183213 & 1.00 & 0.008229 &  \\
base & stretched & 2.845572 & 2.000 & 0.000 & 2.845572 & 0.351423 & 0.25 & 0.602854 & 0.50 & 0.132084 & 1.00 & 0.000304 &  \\
base & logquad & 2.845572 & 1.000 & 0.200 & 2.845572 & 0.351423 & 0.25 & 0.443699 & 0.50 & 0.160791 & 1.00 & 0.011504 &  \\
base & logquad & 2.845572 & 1.000 & 0.500 & 2.845572 & 0.351423 & 0.25 & 0.381199 & 0.50 & 0.087603 & 1.00 & 0.001014 &  \\
base & exp & 4.271140 & 1.000 & 0.000 & 4.271140 & 0.234130 & 0.25 & 0.343769 & 0.50 & 0.118177 & 1.00 & 0.013966 &  \\
base & stretched & 4.271140 & 1.500 & 0.000 & 4.271140 & 0.234130 & 0.25 & 0.331747 & 0.50 & 0.044120 & 1.00 & 0.000147 &  \\
base & stretched & 4.271140 & 2.000 & 0.000 & 4.271140 & 0.234130 & 0.25 & 0.319766 & 0.50 & 0.010455 & 1.00 & 0.000000 &  \\
base & logquad & 4.271140 & 1.000 & 0.200 & 4.271140 & 0.234130 & 0.25 & 0.273674 & 0.50 & 0.047468 & 1.00 & 0.000364 &  \\
base & logquad & 4.271140 & 1.000 & 0.500 & 4.271140 & 0.234130 & 0.25 & 0.194394 & 0.50 & 0.012084 & 1.00 & 0.000002 &  \\
edge & exp & 1.635624 & 1.000 & 0.000 & 1.635624 & 0.611387 & 0.25 & 0.664377 & 0.50 & 0.441396 & 1.00 & 0.194831 & yes \\
edge & stretched & 1.635624 & 1.500 & 0.000 & 1.635624 & 0.611387 & 0.25 & 0.769913 & 0.50 & 0.477318 & 1.00 & 0.123462 &  \\
edge & stretched & 1.635624 & 2.000 & 0.000 & 1.635624 & 0.611387 & 0.25 & 0.846027 & 0.50 & 0.512315 & 1.00 & 0.068889 &  \\
edge & logquad & 1.635624 & 1.000 & 0.200 & 1.635624 & 0.611387 & 0.25 & 0.642527 & 0.50 & 0.386132 & 1.00 & 0.114100 &  \\
edge & logquad & 1.635624 & 1.000 & 0.500 & 1.635624 & 0.611387 & 0.25 & 0.611092 & 0.50 & 0.315935 & 1.00 & 0.051136 &  \\
edge & exp & 2.812226 & 1.000 & 0.000 & 2.812226 & 0.355590 & 0.25 & 0.495070 & 0.50 & 0.245094 & 1.00 & 0.060071 &  \\
edge & stretched & 2.812226 & 1.500 & 0.000 & 2.812226 & 0.355590 & 0.25 & 0.554603 & 0.50 & 0.188744 & 1.00 & 0.008951 &  \\
edge & stretched & 2.812226 & 2.000 & 0.000 & 2.812226 & 0.355590 & 0.25 & 0.610005 & 0.50 & 0.138463 & 1.00 & 0.000368 &  \\
edge & logquad & 2.812226 & 1.000 & 0.200 & 2.812226 & 0.355590 & 0.25 & 0.448470 & 0.50 & 0.165044 & 1.00 & 0.012352 &  \\
edge & logquad & 2.812226 & 1.000 & 0.500 & 2.812226 & 0.355590 & 0.25 & 0.386663 & 0.50 & 0.091201 & 1.00 & 0.001152 &  \\
edge & exp & 3.988828 & 1.000 & 0.000 & 3.988828 & 0.250700 & 0.25 & 0.368908 & 0.50 & 0.136093 & 1.00 & 0.018521 &  \\
edge & stretched & 3.988828 & 1.500 & 0.000 & 3.988828 & 0.250700 & 0.25 & 0.369423 & 0.50 & 0.059810 & 1.00 & 0.000347 &  \\
edge & stretched & 3.988828 & 2.000 & 0.000 & 3.988828 & 0.250700 & 0.25 & 0.369937 & 0.50 & 0.018729 & 1.00 & 0.000000 &  \\
edge & logquad & 3.988828 & 1.000 & 0.200 & 3.988828 & 0.250700 & 0.25 & 0.302374 & 0.50 & 0.061424 & 1.00 & 0.000769 &  \\
edge & logquad & 3.988828 & 1.000 & 0.500 & 3.988828 & 0.250700 & 0.25 & 0.224379 & 0.50 & 0.018625 & 1.00 & 0.000006 &  \\
\bottomrule
\end{tabular}
\caption{Finite-family survival-kernel audit (CAP selection) aligned with the leakage-kernel dictionary.}
\label{tab:survival_finite_family_audit}
\end{table}

\paragraph{Finite-family survival-kernel closure (CAP).} \AuditTag We declare a finite candidate family of survival kernels $P_f(t)$ and CAP-select an element by a deterministic objective. The rate proxy is $\Gamma_f:=\inf\{g>0: P_f(t)\le \exp(-g t)\ \forall t\ge 0\}$; for the declared families it equals the base $\gamma$. Targets are two linewidth proxies $(\Gamma_\tau,\Gamma_{\mathrm{FWHM}})$ obtained from the det-delay carrier. The table reports $\Gamma_f$, $\tau_f=1/\Gamma_f$, and $P_f(t)$ at declared times, and flags the CAP-selected row.
\paragraph{Determinism.} \AuditTag Targets are auto-loaded from \texttt{sections/generated/det\_delay\_linewidth\_rows.tex} (rank=1 per case): base: (Gamma\_tau=4.271140, Gamma\_FWHM=1.420003), edge: (Gamma\_tau=3.988828, Gamma\_FWHM=1.635624); times=[0.25, 0.5, 1.0]. The candidate family and tie-break are finite and explicit.%

\paragraph{Graph-zeta determinant packaging (Ihara/Hashimoto/Bass) and added-edge update (generated audit).}
\AuditTag To close the graph-zeta mother-language into the same certificate registry, we provide a deterministic toy that verifies
the standard determinant identities for Ihara zeta (Hashimoto determinant and Bass determinant form) and reports the added-edge update ratio.
Rows and summary are generated by \texttt{scripts/exp\_ihara\_hashimoto\_added\_edge\_audit.py}.

\begin{table}[t]
\centering
\scriptsize
\setlength{\tabcolsep}{6pt}
\renewcommand{\arraystretch}{1.15}
\begin{tabular}{l r r r r r r}
\toprule
case & $u$ & det(Hashimoto) & det(Bass) & rel\_err & $|E|$ & $|V|$ \\
\midrule
base & 0.050 & 1.000000 & 1.000000 & 0.000000 & 6 & 6 \\
base & 0.100 & 0.999998 & 0.999998 & 0.000000 & 6 & 6 \\
base & 0.150 & 0.999977 & 0.999977 & 0.000000 & 6 & 6 \\
base & 0.200 & 0.999872 & 0.999872 & 0.000000 & 6 & 6 \\
edge & 0.050 & 0.999975 & 0.999975 & 0.000000 & 7 & 6 \\
edge & 0.100 & 0.999598 & 0.999598 & 0.000000 & 7 & 6 \\
edge & 0.150 & 0.997953 & 0.997953 & 0.000000 & 7 & 6 \\
edge & 0.200 & 0.993483 & 0.993483 & 0.000000 & 7 & 6 \\
ratio(edge/base) & 0.050 & 0.999975 & -0.000025 & 0.000000 & 0 & 0 \\
ratio(edge/base) & 0.100 & 0.999600 & -0.000400 & 0.000000 & 0 & 0 \\
ratio(edge/base) & 0.150 & 0.997976 & -0.002026 & 0.000000 & 0 & 0 \\
ratio(edge/base) & 0.200 & 0.993610 & -0.006411 & 0.000000 & 0 & 0 \\
\bottomrule
\end{tabular}
\caption{Ihara/Hashimoto/Bass determinant packaging and added-edge update (toy).}
\label{tab:ihara_hashimoto_added_edge_audit}
\end{table}

\paragraph{Ihara/Hashimoto/Bass determinant packaging and added-edge update (toy).} \AuditTag We compare the Hashimoto determinant $\det(I-uB)$ to the Bass determinant form $(1-u^2)^{|E|-|V|}\det(I-uA+(D-I)u^2)$ on a degree-2 cycle graph $C_6$ and its added-edge (chord) update. We also report the added-edge ratio $\det(I-uB_{\mathrm{edge}})/\det(I-uB_{\mathrm{base}})$ and its logabs increment as a zeta/det update proxy.
\paragraph{Deterministic agreement and update magnitude.} \AuditTag On u-grid [0.05, 0.1, 0.15, 0.2], the maximum relative discrepancy between Hashimoto and Bass determinants is 0.000000 (at case=edge, u=0.200). The maximum observed |Δlog|det|| for the edge/base ratio is 0.006411.%

\paragraph{Hashimoto added-edge update as a $2\times2$ Schur factor (generated audit).}
\AuditTag The Hashimoto determinant package $\det(I-uB)$ lives on the oriented-edge space and therefore changes dimension when an edge is added.
Nevertheless, the update ratio admits a \emph{small} Schur complement factor: with the block split (old oriented edges | new oriented edges),
\[
\frac{\det(I-uB_{\mathrm{new}})}{\det(I-uB_{\mathrm{old}})}=\det(S(u)),
\]
where $S(u)$ is a $k\times k$ Schur complement and $k=2$ in the one-edge update toy.
Rows and summary are generated by \texttt{scripts/exp\_hashimoto\_added\_edge\_schur\_ratio\_audit.py}.

\begin{table}[t]
\centering
\scriptsize
\setlength{\tabcolsep}{6pt}
\renewcommand{\arraystretch}{1.15}
\begin{tabular}{r r r r r r r r}
\toprule
$u$ & det\_old & det\_new & ratio (direct) & ratio (Schur) & err & $k$ & $m_{\mathrm{old}}$ \\
\midrule
0.050 & 1.000000 & 0.999975 & 0.999975 & 0.999975 & 0.000000 & 2 & 12 \\
0.100 & 0.999998 & 0.999598 & 0.999600 & 0.999600 & 0.000000 & 2 & 12 \\
0.150 & 0.999977 & 0.997953 & 0.997976 & 0.997976 & 0.000000 & 2 & 12 \\
0.200 & 0.999872 & 0.993483 & 0.993610 & 0.993610 & 0.000000 & 2 & 12 \\
\bottomrule
\end{tabular}
\caption{Hashimoto determinant added-edge ratio via a small Schur complement factor (toy).}
\label{tab:hashimoto_added_edge_schur_audit}
\end{table}

\paragraph{Added-edge update ratio via Schur complement (Hashimoto determinant).} \AuditTag Adding one undirected edge increases the oriented-edge space, so the Hashimoto matrix $B$ grows in dimension. With the block split (old oriented edges | new oriented edges), the block determinant identity yields $\det(I-uB_{\mathrm{new}})=\det(I-uB_{\mathrm{old}})\det(S(u))$ where $S(u)$ is a $k\times k$ Schur complement with $k=2$ in this toy. We compare the direct ratio to $\det(S(u))$ on a declared $u$ grid.
\paragraph{Deterministic agreement.} \AuditTag On u-grid [0.05, 0.1, 0.15, 0.2], the maximum absolute discrepancy is 0.000000 (at u=0.050).%
